

\chapter{Arquitetura do Sistema}
\label{chapter:architecture}

\begin{introduction}
Este capítulo apresenta a visão geral da arquitetura do sistema, descrevendo os subsistemas mecânico, elétrico e de software que o compõem.
\end{introduction}



\section{Visão Geral da Arquitetura}
\label{sec:visao-geral}

O sistema divide-se em três subsistemas interligados. O subsistema mecânico trata da movimentação física da plataforma; o subsistema elétrico e eletrónico controla os atuadores e gere a energia; o subsistema de software e firmware abarca o pré-processamento dos dados sísmicos e a lógica de controlo em tempo real.


\section{Subsistema Mecânico}
\label{sec:subsistema-mecanico}

O subsistema mecânico é responsável por converter o movimento rotativo dos motores em deslocamento linear controlado da plataforma, em dois eixos horizontais ortogonais. A estrutura é composta por dois elementos principais: uma base fixa, fabricada por impressão 3D, que suporta os quatro motores e serve de fundação rígida do sistema; e uma plataforma móvel única, também impressa em 3D, sobre a qual são colocados os modelos estruturais a ensaiar.

A ligação entre a base e a plataforma é assegurada por quatro colunas de suporte, dotadas de rótulas nas duas extremidades (topo e base). Esta configuração é o elemento central do design bi-axial: ao contrário de soluções com guias lineares rígidas por eixo, as rótulas permitem que as colunas se inclinem livremente em qualquer direção, acomodando o movimento simultâneo da plataforma nos eixos~X e Y sem introduzir constrangimentos cinemáticos entre eles.

Cada motor está posicionado próximo do centro da base, orientado radialmente para o exterior (posições N, S, E e W), e acoplado a um fuso trapezoidal com \textit{lead} de 8\,mm por revolução. Em cada extremidade da plataforma está fixada uma porca trapezoidal correspondente: quando o motor roda, o fuso avança ou recua na porca, empurrando ou puxando a plataforma nessa direção. Os motores N e S controlam o eixo~Y (componente Norte-Sul do sismo), rodando em sentidos opostos de modo a que ambos apliquem força na mesma direção sobre a plataforma. O mesmo princípio aplica-se aos motores E e W para o eixo~X (componente Este-Oeste). Os quatro motores funcionam em simultâneo, pelo que o movimento bi-axial resulta da composição vetorial das forças aplicadas pelos dois pares de motores em cada instante.


\section{Subsistema Elétrico e Eletrónico}
\label{sec:subsistema-eletrico}

O subsistema elétrico e eletrónico centraliza-se no microcontrolador ESP32, que funciona como unidade central de controlo e processamento. A CNC Shield, encaixada diretamente sobre o ESP32, providencia a interface física e elétrica entre o microcontrolador e os quatro \textit{drivers} DRV8825, um por motor. Esta arquitetura modular reduz a complexidade da cablagem e facilita a substituição de componentes em caso de avaria.

A alimentação do sistema é bipartida: o ESP32 é alimentado pela porta USB durante o desenvolvimento e operação normal, enquanto os motores recebem energia de uma fonte DC externa dedicada, com tensão e corrente adequadas aos motores NEMA17 utilizados. Separar eletricamente o domínio lógico (ESP32 e lógica dos \textit{drivers}) do domínio de potência (enrolamentos dos motores) evita perturbações nos sinais de controlo que poderiam causar comportamentos erráticos no firmware.

O pino Enable (EN) dos \textit{drivers}, partilhado entre todos os quatro \textit{drivers} através da CNC Shield e ligado ao pino 13 do ESP32, é ativado (nível lógico baixo) antes do início de cada sequência de movimento e desativado (nível lógico alto) no final, permitindo o arrefecimento dos motores e a redução do consumo energético durante os períodos de inatividade.


\section{Subsistema de Software e Firmware}
\label{sec:subsistema-software}

O subsistema de software e firmware divide-se em dois componentes com ciclos de vida distintos. O script de pré-processamento em Python opera em modo \textit{offline}, executado uma única vez na fase de preparação do sistema, e converte os dados sísmicos de deslocamento em arrays de velocidades de motor, prontos a ser embebidos no firmware. O firmware em C++ para o ESP32, por outro lado, corre em tempo real durante os ensaios, percorrendo os arrays pré-calculados e controlando os motores a uma cadência de 50\,Hz.

A biblioteca AccelStepper é utilizada no firmware para o controlo de velocidade instantânea dos motores, em modo \texttt{runSpeed()}, que aplica a velocidade definida sem qualquer rampa de aceleração ou desaceleração. Esta escolha justifica-se porque os dados sísmicos codificam implicitamente a variação de velocidade entre janelas consecutivas, pelo que qualquer aceleração adicional imposta pela biblioteca seria redundante e potencialmente distorciva do perfil sísmico real.

A comunicação série, à taxa de 115\,200\,baud, serve como canal de telemetria, permitindo a monitorização em tempo real das velocidades instantâneas dos quatro motores durante a execução de uma sequência sísmica. O formato de saída CSV possibilita a captura e análise posterior dos dados com ferramentas como Python ou Excel.
